% Sin número: se conserva la numeración de las secciones del documento.
\section*{Componentes de la plantilla de Blueprint}
\addcontentsline{toc}{section}{Componentes de la plantilla de Blueprint}

El concepto de Blueprint describe qué es y para qué sirve; la plantilla lo traduce en un conjunto de
bloques que se desarrollan en un orden determinado. Cada bloque responde a una pregunta distinta y
alimenta al siguiente: qué se quiere lograr y hasta dónde llega el diseño, qué contexto y obligaciones
lo condicionan, de dónde se parte y qué falta, qué amenazas y activos justifican los controles, cómo se
pasa al estado objetivo y cómo se comprueba el avance. La Tabla~\ref{tab:plantilla_blueprint} resume
los seis bloques de la plantilla.

\begin{table}[H]
	\caption{Componentes de la plantilla de Blueprint}
	\label{tab:plantilla_blueprint}
	\centering
	\footnotesize
	\renewcommand{\arraystretch}{1.15}
	\begin{tabularx}{\textwidth}{@{}>{\raggedright\arraybackslash}p{0.22\textwidth} Y >{\centering\arraybackslash}p{0.06\textwidth}@{}}
		\toprule
		\textbf{Bloque} & \textbf{Qué define} & \textbf{Sec.} \\
		\midrule
		\textbf{Visión, principios y alcance}
		& Fija la razón de ser del Blueprint, las fronteras del diseño, las reglas no negociables y los interesados que deben validarlo.
		& 1
		\\
		\midrule
		\textbf{Contexto de negocio, objetivos y cumplimiento}
		& Relaciona las decisiones de seguridad con los objetivos del negocio, la normativa aplicable y el nivel de riesgo que la entidad acepta.
		& 2
		\\
		\midrule
		\textbf{Evaluación del estado actual y análisis de brechas}
		& Describe la arquitectura base, la evalúa en efectividad, madurez y eficiencia, y compara el estado actual con el objetivo.
		& 3
		\\
		\midrule
		\textbf{Modelo de amenazas, activos y evaluación de riesgos}
		& Identifica los activos críticos, caracteriza la superficie de ataque y valora el riesgo inherente y residual.
		& 4
		\\
		\midrule
		\textbf{Arquitectura de transición y hoja de ruta}
		& Define los estados intermedios, el cronograma, los entregables y los recursos para pasar del estado actual al objetivo.
		& 5
		\\
		\midrule
		\textbf{Métricas}
		& Establece los indicadores de desempeño y de riesgo, con metas y umbrales, para verificar el cumplimiento de los objetivos.
		& 6
		\\
		\bottomrule
	\end{tabularx}

	\vspace{0.4em}
	\parbox{\textwidth}{\footnotesize \textit{Nota.} Elaboración propia.}
\end{table}

La plantilla incluye dos componentes complementarios que se integran en los bloques anteriores: las
políticas de seguridad de O-ESA y el perfil de atributos de negocio de SABSA. Las políticas se reflejan
en la matriz SABSA y en el esquema PDP/PEP de la Sección~5; los atributos de negocio se concretan en
los KPI y KRI de la Sección~6.

Los bloques se encadenan: la visión y el alcance definen el perímetro; el contexto y el cumplimiento
fijan los criterios; el diagnóstico y las brechas establecen el punto de partida; el modelo de
amenazas y riesgos justifica los controles; la transición y la hoja de ruta ordenan la implementación;
y las métricas permiten comprobar si el estado objetivo se alcanzó. Cada bloque conserva la referencia
al anterior, de modo que una decisión de diseño pueda rastrearse hasta el objetivo, el riesgo o la
obligación que la origina.

Las secciones siguientes desarrollan estos bloques en ese mismo orden, comenzando por la visión, los
principios y el alcance de la arquitectura.
